% Detailed manuscript outline based on the current experiment code and saved results.
% This is an outline, not a submission-ready manuscript. Replace bracketed
% author/venue fields, add verified citations, and report final values from the
% aggregate outputs before drafting prose.
\documentclass[pdflatex,sn-mathphys-num]{sn-jnl}
\usepackage{amsmath}

\title[Thresholds, resampling, and credit-default evaluation]{
When Does Resampling Help? A Multi-Dataset Evaluation of Credit-Default Models,
Decision Thresholds, and Reliability Diagnostics
}
\author{[Author names and affiliations to be added]}
\date{}

\abstract{\textbf{Outline.} Motivate the practical problem of comparing credit-default models under class imbalance without allowing test data to guide preprocessing, tuning, or threshold choice. State the study design: six classifiers, ten random seeds, and three datasets, including a row-chronological Lending Club experiment. Preview the central empirical result: validation-tuned thresholds improve F1 over a fixed 0.50 threshold, while the benefit of synthetic resampling varies by dataset. Report final aggregate estimates and uncertainty only after reconciling all tables and figures. Describe robustness, calibration, fairness, and explanation diagnostics as supporting evaluations; do not present the current descriptive composite score as a validated deployment rule.}
\keywords{Credit scoring, default prediction, class imbalance, decision threshold, model evaluation, explainable AI, fairness}

\begin{document}
\maketitle

% -----------------------------------------------------------------------------
% Central argument and scope
% -----------------------------------------------------------------------------
\noindent\textbf{Proposed central contribution.} A leakage-conscious, repeated multi-dataset comparison shows that decision-threshold selection and resampling are distinct interventions: validation-tuned thresholds improve minority-class operating metrics in the evaluated datasets, while oversampling is not consistently beneficial and its value differs across datasets. The paper contributes comparative evidence and a reporting protocol, not a new classifier or a universal credit-scoring standard.

\noindent\textbf{Primary research question.} Across credit-default datasets, how do model ranking, operating-threshold selection, and training-set resampling jointly affect predictive performance, calibration, and minority-class detection when model development is separated from final test evaluation?

\noindent\textbf{Secondary questions.}
\begin{enumerate}
  \item How stable are model comparisons across random seeds and across random versus temporal evaluation?
  \item How do input perturbations affect predictive performance and individual decisions?
  \item How do SHAP and LIME compare on attribution concentration, local fidelity, agreement, and stability under the implemented diagnostics?
  \item What utility--fairness trade-offs appear for datasets with demographic or group attributes?
\end{enumerate}

\noindent\textbf{Evidence boundary.} UCI and South German runs use randomized train/validation/test splits; Lending Club uses a row-chronological split, with same-month loans potentially crossing a split boundary. Results are repeated across ten seeds, which vary both the split and stochastic training. Avoid describing those repeats as ten evaluations on one fixed test set.

\section{Introduction}
\textit{Target length: about 1200 words.}

\subsection{Credit-default prediction as a decision problem}
\textbf{Purpose.} Explain why a model score is not itself a lending decision: a threshold converts probability estimates into actions, and the appropriate operating point depends on error costs and institutional objectives.

\textbf{Content.} Introduce class imbalance and the common use of resampling. Distinguish ranking quality (for example, AUROC/AUPRC), probability quality (Brier score and calibration), and threshold-dependent measures (F1, recall, specificity, expected utility). Use a brief intuitive explanation of why changing a threshold changes recall and false-positive burden without necessarily changing AUROC.

\textbf{Evidence needed.} Credit-scoring literature on model evaluation, class imbalance, threshold choice, calibration, and the cost of false approvals/false declines. Add verified references; do not infer literature claims from the code.

\subsection{Why a leakage-conscious comparison is needed}
Describe the risks of fitting preprocessing, selecting thresholds, or selecting resampling strategies using held-out test data. State the intended separation: train-only preprocessing and resampling; validation for model/threshold/policy choices; held-out test for final descriptive evaluation. Explain that the design is intended to make comparisons interpretable, not to certify the models as deployment-ready.

\subsection{Research gap and contribution}
Position the work as an empirical comparison, not as a new algorithm. Explain the gap the paper addresses: the practical tendency to discuss oversampling without isolating its effects from threshold selection, calibration, and data-split variability. Compare with prior credit-scoring benchmarks and recent work on integrated AI-risk metrics and formal XAI comparison criteria. Relevant anchors include the SAFE AI metrics paper (DOI: 10.1016/j.mlwa.2025.100821) and the explainable-AI assessment framework (DOI: 10.1007/s10462-025-11141-w); add accurate bibliographic records and discuss differences rather than claiming equivalence.

\textbf{Contribution list.}
\begin{enumerate}
  \item A repeated, three-dataset comparison of six classifiers with separate training, validation, and test roles.
  \item A controlled empirical comparison of no resampling and four over-sampling strategies, alongside fixed, validation-tuned, and validation-fitted calibrated decision rules.
  \item Supporting diagnostics for robustness, calibration, explanation quality, and fairness, reported with their scope and limitations.
\end{enumerate}

\subsection{Research questions and paper roadmap}
List the primary and secondary questions above in compact form, then give one sentence per remaining section. State explicitly that findings are dataset-dependent and that statistical significance is not treated as practical importance.

\section{Related Work and Conceptual Positioning}
\textit{Target length: about 1050 words.}

\subsection{Credit scoring and comparative classifiers}
Review conventional scorecards/logistic regression, tree ensembles, and neural tabular models. Summarize common evaluation practices and distinguish this study's three-dataset, multi-seed design from prior single-dataset studies. Identify the closest benchmarks and compare dataset, split, metrics, models, and explainability scope in a compact related-work table.

\subsection{Class imbalance and decision thresholds}
Review resampling methods (SMOTE, Borderline-SMOTE, SVM-SMOTE, ADASYN) and alternatives based on decision thresholds and costs. Explain mathematically that a probability threshold changes predicted labels while leaving the ranking of probabilities, and thus ROC-AUC, unchanged for a fixed fitted model. Clarify that resampling changes the fitted model and probability distribution, so it can affect both ranking and calibration.

\subsection{Calibration, robustness, and reliability}
Define probability calibration and distinguish it from discrimination. Review perturbation-based robustness and integrated metrics, identifying which SAFE AI dimensions are and are not examined here. The present code measures missingness, correlated missingness, numeric noise, and covariate shift; it does not implement a poisoning/security evaluation.

\subsection{Explanation comparison and fairness trade-offs}
Review global and local explanations, parsimony/concentration, fidelity, and stability. Explain that feature attribution is not causal evidence or a regulatory justification by itself. Review accuracy/fairness trade-off work and specify group attributes and metrics relevant to the available datasets. Do not equate a threshold post-processor or one gap metric with legal compliance.

\section{Study Design and Methods}
\subsection{Research framework}
The study uses a repeated, multi-dataset benchmarking design. For each dataset and seed, the pipeline loads and cleans observations, creates training, validation, and test partitions, fits preprocessing on the training partition, compares resampling strategies, tunes classifiers, selects a model and operating threshold using validation data, and evaluates the frozen choices on the held-out test partition. Additional analyses assess probability calibration, input perturbations, feature-attribution diagnostics, and group disparities. The test partition is not used to select a model, sampler, calibration map, threshold, or fairness policy. The primary estimand is the held-out performance of each fitted model under the validation-selected decision threshold; seed-to-seed variation includes both partition variation and stochastic model fitting.

\subsection{Datasets and prediction targets}
Three tabular datasets were used. The UCI Default of Credit Card Clients file contains 30,000 rows and 25 columns, including an identifier, 23 recorded predictors, and the target \texttt{default payment next month}. The identifier is discarded. The positive class is target value 1 (default); the local file contains 6,636 positive and 23,364 negative observations (22.12\% positive). The predictor set includes credit limit, sex, education, marital status, age, six repayment-status variables, six statement balances, and six payment amounts. The data file is bundled as \texttt{data/uci\_credit.csv}. Its citation and exact source-file release/version should be supplied in the final reference list [MISSING INFORMATION].

The corrected South German Credit file contains 1,000 observations, 20 predictors, and one credit-risk label. The pipeline maps the source's numeric coding to a binary outcome in which bad credit is 1 and good credit is 0; the bundled file encodes 700 good and 300 bad cases, giving 30\% positive prevalence. The predictors are the 20 named fields in the supplied UCI file, including status, duration, credit history, purpose, amount, savings, employment duration, personal status/sex, age, housing, telephone, and foreign-worker status. The pipeline accepts \texttt{data/SouthGermanCredit.asc} or two alternative local CSV names and, if no local file is available, downloads the corrected dataset from UCI. Add the exact source citation/version used [MISSING INFORMATION].

The bundled Lending Club sample has 50,000 rows and 18 columns before preprocessing. The temporal key spans July 2007 through December 2018. The data loader removes the row identifier, issue-date text, and loan-status field; converts percentage strings in \texttt{int\_rate} and \texttt{revol\_util} to numeric values; constructs \texttt{issue\_time} from the issue date when needed; and drops rows with missing values. The final model matrix contains 13 predictors: loan amount, funded amount, interest rate, installment, annual income, debt-to-income ratio, delinquency count, recent inquiries, open accounts, public-record count, revolving balance, revolving utilization, and total accounts. The packaged target is \texttt{default\_payment\_next\_month}, with 9,901 positive (default) and 40,099 negative observations. The exact rule used to create this precomputed target, the original Lending Club source/version, and sample selection procedure are not recorded in the repository [MISSING INFORMATION]; the run manifest also records no checksum for this file.

For UCI, group analyses use \texttt{SEX}, \texttt{EDUCATION}, and \texttt{MARRIAGE}; for South German they use \texttt{personal\_status\_sex}, \texttt{age}, and \texttt{foreign\_worker}. These variables remain predictors in the fitted models. Lending Club group fairness is omitted because no supported group fields are available in the packaged sample. These group fields identify observed dataset attributes only; their use does not define protected classes for legal or policy evaluation.

\subsection{Data preprocessing and partitioning}
The ten run seeds are 42, 99, 123, 326, 456, 515, 689, 777, 872, and 999. UCI and South German use two successive stratified random splits: 20\% of all observations are assigned to test, then 10\% of the full dataset is selected from the remaining 80\% for validation. This gives 70\%/10\%/20\% train/validation/test partitions. The Lending Club runs sort rows by \texttt{issue\_time} and take the earliest 70\% for training, the next 10\% for validation, and the latest 20\% for test. This is a row-chronological split, not a month-grouped split; records sharing an issue month may fall on both sides of a boundary. The temporal key is removed from model predictors. Partition indices are saved for each run.

UCI cleaning drops \texttt{ID}, renames the target consistently, maps education codes 0, 5, and 6 to 4, and maps marriage code 0 to 3. No imputation is applied to UCI or South German. The Lending Club loader drops incomplete rows rather than imputing. The South German categorical fields are retained in their numeric source coding; no one-hot encoding is applied. All numeric predictors, including integer-coded categories, are standardized using means and standard deviations fitted on the training partition; the fitted transform is then applied without refitting to validation and test data.

For the UCI data, eight engineered predictors are added when the required source fields are present: mean utilization (mean statement balance divided by credit limit), mean payment-to-positive-bill ratio, mean and maximum nonnegative repayment delay, linear slopes of six-month bill and payment histories, credit-limit-to-mean-bill ratio, and an indicator that the most recent delay exceeds the mean of earlier delays. Negative repayment-status codes are clipped to zero for delay summaries; payment ratios with nonpositive bills are set to zero. No domain feature block is added to the other datasets. Numeric predictors are winsorized at the training partition's 1st and 99th percentiles, with those bounds reused unchanged for validation and test. The code calculates the engineered fields and winsorization bounds after the outer split, so held-out rows do not determine the bounds. It does not fit preprocessing separately inside the randomized-search folds.

Each seed changes the random partitions for UCI and South German, as well as sampler and model random states. For Lending Club the partition order is fixed by time, although seed-specific model and analysis operations remain. Consequently, the ten runs are repeated protocol executions, not ten independent datasets or ten tests on one identical held-out sample.

\subsection{Classifiers, baselines, and hyperparameter selection}
The six classifier families are XGBoost, LightGBM, CatBoost, Random Forest, Logistic Regression, and TabNet. Logistic Regression is the linear scorecard-style baseline; Random Forest is a bagged-tree baseline; the three boosting methods represent gradient-boosted trees; and TabNet is the neural tabular baseline. The source config sets class weights to \texttt{balanced} for Logistic Regression and Random Forest, while boosting models and TabNet do not use this setting. Since models therefore do not share identical imbalance handling or optimization procedures, the experiment is a benchmark of the configured pipelines rather than a controlled comparison under equal computational budgets.

For XGBoost, LightGBM, CatBoost, Random Forest, and Logistic Regression, the main runs use \texttt{RandomizedSearchCV} with 12 sampled configurations per model, three-fold shuffled stratified CV, ROC-AUC scoring, and seed-controlled fold and parameter sampling. The search is run on the outer training matrix after the resampling strategy has been applied. Search ranges are defined in \texttt{config.py}: XGBoost varies tree count (300, 500, 800), depth (3, 4, 5, 6, 8), learning rate (0.01--0.10), row/column subsampling, gamma, and L1/L2 regularization; LightGBM varies tree count, depth, leaves (15--127), learning rate, subsampling, regularization, and minimum child samples (10--80); CatBoost varies iterations (300, 500, 800), depth (4, 5, 6, 8), learning rate, and L2 regularization; Random Forest varies tree count (200, 300, 500), depth, split/leaf minima, and feature-sampling mode; Logistic Regression varies inverse regularization $C$ over $[0.01,10]$. Exact discrete candidate values are preserved in the cited configuration file and should accompany a reproducibility archive. If search fails, the implementation fits the fixed model configuration and records a fallback status. The selected configuration is refit on the complete outer training partition. XGBoost, LightGBM, and CatBoost wrappers may use validation-based early stopping according to their configured settings.

TabNet is trained once per seed with fixed settings (64-dimensional decision and attention representations, five steps, Adam optimizer, learning rate 0.02, maximum 200 epochs, patience 15, batch size 1,024, and virtual batch size 128); it is not included in randomized hyperparameter search. A failed or skipped fit is recorded in the run output. The fixed settings and actual status should be reported from each run manifest before final submission [MISSING INFORMATION: per-seed TabNet fit-status summary].

\subsection{Resampling, operating thresholds, and calibration}
Five training strategies are compared: no resampling, SMOTE, Borderline-SMOTE (borderline-1), SVM-SMOTE, and ADASYN. The samplers use five nearest neighbors and the run seed where supported. Each sampler is fitted only on the outer training data. To choose the main-run strategy, the pipeline fits each resampled training set with one available tree model (LightGBM first, then XGBoost, CatBoost, or Random Forest as fallback), scores its validation AUROC, and selects the highest-AUROC strategy. The selected resampled matrix is then used to train each main classifier. Validation and test rows are not synthetically resampled.

This order has an important reproducibility qualification: resampling precedes the three-fold randomized search. Synthetic observations therefore enter the search matrix before folds are formed, and related original/synthetic observations may be split across internal folds. This can make internal CV scores for resampled conditions optimistic. The outer validation and test partitions remain separate from fitting samplers, but the internal resampled CV procedure is not sampler-within-fold cross-validation; the results must not be described as fully leakage-free model selection.

For each fitted model, the validation set determines the operating threshold. Candidate thresholds range from 0.05 to 0.94 in increments of 0.01; the threshold with the largest validation F1 is selected (the first threshold is retained in a tie). The per-seed champion is chosen by validation AUROC among fitted models. Its test predictions are then evaluated at its validation-selected threshold. Thus, validation F1 selects the decision rule, while validation AUROC selects the champion; neither choice uses test metrics.

A separate controlled ablation refits the validation-selected champion family, using its fixed/default configuration, on each of the five resampled training sets. It compares a fixed 0.50 threshold, a threshold maximizing validation F1, and isotonic calibration fitted to validation predictions followed by an F1-maximizing threshold on calibrated validation probabilities. The calibration mapping and threshold are frozen before evaluating the test partition. This controlled ablation is distinct from the main run, which uses tuned model configurations and the validation-selected sampler. A raw-feature versus engineered-feature ablation is run only for UCI, with the same split and champion family but without a second hyperparameter search.

\subsection{Evaluation metrics}
For binary labels $y_i\in\{0,1\}$ and estimated default probabilities $p_i$, AUROC measures the probability that a randomly selected positive receives a higher score than a randomly selected negative. AUPRC is average precision, the weighted average of precision at the recall increments. These summarize ranking and do not depend on a single threshold. At threshold $t$, predictions are $\hat y_i(t)=\mathbf{1}(p_i\ge t)$. Precision is $TP/(TP+FP)$, recall/sensitivity is $TP/(TP+FN)$, specificity is $TN/(TN+FP)$, and $F_1=2PR/(P+R)$. G-Mean is $\sqrt{\mathrm{recall}\,\mathrm{specificity}}$; Matthews correlation coefficient (MCC) summarizes the four confusion-matrix cells. The Brier score, $n^{-1}\sum_i(p_i-y_i)^2$, measures squared probability error.

Calibration diagnostics use ten equal-width probability bins. Expected calibration error (ECE) is the sample-size-weighted mean absolute difference between each bin's mean prediction and observed event rate; maximum calibration error (MCE) is the largest such difference. Calibration intercept and slope are estimated from logistic regression of outcomes on the logit of predicted probabilities; calibration-in-the-large estimates the intercept with slope fixed at one. Isotonic regression is used only in the controlled calibration ablation. The reported metrics address different properties: ranking, threshold decisions, and probability accuracy are not interchangeable.

The fairness/utility curves evaluate thresholds 0.20 to 0.60 in 0.05 increments. Per-observation utility is $(TP-FP-5FN)/n$, corresponding to true-positive benefit 1, true-negative benefit 0, false-positive cost 1, and false-negative cost 5. For each observed group attribute, selection-rate, true-positive-rate, and false-positive-rate gaps are the maximum minus minimum rates across groups. A separate equal-opportunity threshold post-processor searches group thresholds on validation data over candidate target TPRs from 0 to 1; it aims for a validation TPR gap no larger than 0.10 and selects the feasible policy with highest validation utility. The frozen policy is evaluated on test data. These definitions are descriptive and depend on the chosen group fields and utility costs.

\subsection{Robustness and explanation diagnostics}
Robustness is assessed by perturbing test predictors at levels 0.05, 0.10, and 0.20. The four implemented perturbations are: independent numeric-cell missingness replaced with the training median; correlated block missingness in which repayment/billing/utilization fields or other numeric fields are jointly replaced by training medians; Gaussian numeric noise with standard deviation equal to the level times the training standard deviation; and covariate shift that adds the level to monetary, utilization, income, debt-to-income, or FICO-like predictors. The UCI fields \texttt{SEX}, \texttt{EDUCATION}, and \texttt{MARRIAGE} are excluded from perturbation. Other integer-coded categorical or group fields are treated as numeric when perturbed; no clipping is applied after noise or shift. Each perturbed set is scored with the original labels and the validation-selected threshold. Diagnostics include AUROC retention relative to the unperturbed test AUROC, Brier retention defined as $1-(Brier_{perturbed}-Brier_{base})/Brier_{base}$, mean absolute probability shift, and the fraction of thresholded decisions changed. These are specified synthetic corruption tests, not external temporal validation or adversarial robustness tests.

Explanation analyses are performed for the validation-selected champion, not for all six model families. The revision diagnostics use up to 100 test instances per seed. SHAP computes feature attributions for the fitted model; LIME fits a local linear surrogate using 1,000 perturbed samples per instance, discretizes continuous features, and returns up to 15 features. The analysis records local LIME $R^2$, absolute-attribution rank correlation between SHAP and LIME on features present in the LIME explanation (only when at least three shared features exist), attribution parsimony (Gini concentration and number of features accounting for 80\% of absolute attribution mass), and LIME stability across three repeated explanations of the same observation/model. Stability is summarized by pairwise Spearman rank correlation and top-10 Jaccard overlap. The repeated explanations do not refit the model. Fidelity diagnostics replace the most important 20\% of features by training medians: comprehensiveness is the absolute prediction change after deleting these features, while sufficiency is the absolute difference when only these features are retained and remaining values use the median reference. Attributions describe model behavior; they are not causal effects or evidence that a decision is justified.

\subsection{Statistical analysis and uncertainty}
The main tables summarize metrics as means and standard deviations across the ten seed runs. Within each seed, pairwise AUROC differences are tested with DeLong's test. Paired classification correctness is compared with continuity-corrected McNemar tests, using the validation-selected thresholds. Holm step-down adjustment is applied separately to the AUROC and McNemar pairwise-test families within each run. When revision analyses are enabled, percentile intervals are estimated from 2,000 paired row-bootstrap samples on the held-out test set; each bootstrap resamples rows jointly across models and reports 2.5th and 97.5th percentiles for AUROC, AUPRC, Brier, F1, and MCC, as well as paired model differences. Group-policy differences use the same row-resampling principle with the policy fixed. These intervals condition on each run's fitted models and test sample. Across-seed standard deviations instead reflect changes in splits and training randomness. The code does not combine the ten seed-level hypothesis tests into a single pooled inferential test.

\subsection{Composite reliability diagnostics}
The code produces two distinct reliability summaries. During validation-time candidate screening, four components are combined by their geometric mean after clipping to $[0,1]$: AUROC; Brier skill, calculated as $1-Brier/Brier_{base}$ using the validation prevalence; mean AUROC retention under validation perturbations; and $1-$maximum group-rate gap at the threshold grid value nearest the model's validation-selected threshold. A candidate is feasible only when AUROC retention is at least 0.90 and the maximum gap is at most 0.10. The pipeline ranks feasible candidates first and otherwise records an infeasible fallback. This validation screen does not include explainability.

A separate test-set audit score is computed only for the validation-selected champion. Its four clipped components are predictive quality, defined as $\sqrt{AUROC\times BrierSkill}$; mean test AUROC retention across perturbations; an explanation aggregate equal to the mean of normalized LIME rank stability, SHAP parsimony, comprehensiveness, and LIME local $R^2$; and one minus the largest available TPR/FPR gap. The geometric mean is reported only when group fairness is available. Component weights, normalization, and thresholds are exploratory implementation choices. The score is not calculated for every model in a comparable common validation framework and is not a deployment decision rule.

\subsection{Reproducibility record and missing details}
The experiment was invoked through \texttt{run\_all\_full.sh} and \texttt{pipeline.py} for each dataset, using the ten seeds above, 12 search iterations, revision analyses, up to 100 explanation instances, and 2,000 bootstrap replicates. Split indices, per-run results, source-code fingerprints, and some dataset checksums are written under \texttt{results/outputs/<dataset>/runs/}. The manuscript's reproducibility package should include the source code, configuration, input data or authorized retrieval instructions, split indices, and per-seed model/search records.

The repository does not contain an exact environment lockfile; \texttt{requirements.txt} specifies lower bounds rather than the installed package versions [MISSING INFORMATION]. The Lending Club input checksum and the provenance, inclusion criteria, and target-construction rule for the prepackaged 50,000-row sample are absent [MISSING INFORMATION]. Some manifests store machine-specific absolute output paths, and therefore need portable relative-path copies for archival release [MISSING INFORMATION]. These details are needed to reproduce the precise data artifact and software environment.

\section{Results}
\textit{Results below summarize the saved ten-seed runs. Means and standard deviations describe variation across seeds; they are not confidence intervals over a single fixed test set.}

\subsection{Dataset composition and model performance}
Each dataset was evaluated on the same partition sizes across the ten seeds: UCI used 21,000 training, 3,000 validation, and 6,000 test observations; South German used 700, 100, and 200; and the packaged Lending Club sample used 35,000, 5,000, and 10,000. Test-set default prevalence was 22.12\%, 30.00\%, and 21.58\%, respectively. Model metrics below use the pipeline's validation-selected resampling strategy and validation-F1 threshold for each run.

Table~\ref{tab:all-seed-models} reports the mean test AUROC, AUPRC, F1, and Brier score for each model across ten seeds. On UCI, the leading mean AUROC was LightGBM (0.7879, SD 0.0050), closely followed by CatBoost (0.7877), XGBoost (0.7871), and Random Forest (0.7860). LightGBM also had the highest mean AUPRC (0.5694) and lowest mean Brier score (0.1325); CatBoost had the highest mean F1 (0.5500). The validation-selected champion varied across seeds: XGBoost won four runs, LightGBM three, and CatBoost three. Thus, UCI supports a narrow advantage over Logistic Regression, not a clear winner among the leading ensemble models.

On South German, Random Forest had the highest mean AUROC (0.7901, SD 0.0216); Logistic Regression was close (0.7835), and had the highest mean AUPRC (0.6342). Random Forest also had the highest mean F1 (0.5929) and lowest Brier score (0.1684). Performance varied more across seeds than on UCI: the validation-selected champion's test AUROC ranged from 0.7144 to 0.8048, and champion identity varied across five model families. TabNet was notably weaker on average (AUROC 0.6566, SD 0.0450; Brier 0.2904). The results do not support a general claim that gradient boosting leads on this dataset.

On temporal Lending Club, CatBoost had the highest mean AUROC (0.6754, SD 0.0008) and F1 (0.4228); XGBoost had the highest AUPRC (0.3439), only 0.0003 above CatBoost. Differences among the leading models were small. Champion identity varied across seeds (CatBoost five, XGBoost four, LightGBM one), while champion test AUROC ranged from 0.6745 to 0.6762. The low AUROC indicates modest discrimination under this row-chronological split, despite stable estimates across seeds.

Holm-adjusted DeLong tests were conducted separately within each seed. On UCI, XGBoost, LightGBM, CatBoost, and Random Forest each exceeded Logistic Regression in AUROC in all ten seed-level tests after adjustment; pairwise differences among the leading ensembles were rarely significant. No model pair was significant in all ten South German tests, and no Lending Club pair remained significant in all ten after adjustment. These are repeated within-split tests, not a pooled test across seeds. Table values therefore support dataset-dependent rankings and narrow ensemble differences, not universal superiority.

\begin{table}[ht]
\centering\small
\caption{Held-out performance across ten seeds. Values are mean test metrics; AUROC standard deviation across seeds is shown in parentheses.}
\label{tab:all-seed-models}
\begin{tabular}{llcccc}
\hline
Dataset & Model & AUROC (SD) & AUPRC & F1 & Brier \\
\hline
UCI & CatBoost & 0.7877 (0.0062) & 0.5686 & 0.5500 & 0.1384 \\
UCI & LightGBM & 0.7879 (0.0050) & 0.5694 & 0.5495 & 0.1325 \\
UCI & Logistic Reg. & 0.7665 (0.0059) & 0.5208 & 0.5330 & 0.1876 \\
UCI & Random Forest & 0.7860 (0.0058) & 0.5677 & 0.5456 & 0.1746 \\
UCI & TabNet & 0.7780 (0.0057) & 0.5470 & 0.5456 & 0.1353 \\
UCI & XGBoost & 0.7871 (0.0066) & 0.5677 & 0.5483 & 0.1351 \\
\hline
South German & CatBoost & 0.7752 (0.0220) & 0.6127 & 0.5694 & 0.1971 \\
South German & LightGBM & 0.7760 (0.0208) & 0.6213 & 0.5845 & 0.1717 \\
South German & Logistic Reg. & 0.7835 (0.0244) & 0.6342 & 0.5899 & 0.1843 \\
South German & Random Forest & 0.7901 (0.0216) & 0.6334 & 0.5929 & 0.1684 \\
South German & TabNet & 0.6566 (0.0450) & 0.4484 & 0.4835 & 0.2904 \\
South German & XGBoost & 0.7764 (0.0276) & 0.6067 & 0.5634 & 0.1965 \\
\hline
Lending Club & CatBoost & 0.6754 (0.0008) & 0.3436 & 0.4228 & 0.1603 \\
Lending Club & LightGBM & 0.6740 (0.0010) & 0.3431 & 0.4167 & 0.1599 \\
Lending Club & Logistic Reg. & 0.6704 (0.0001) & 0.3386 & 0.4195 & 0.2353 \\
Lending Club & Random Forest & 0.6744 (0.0014) & 0.3381 & 0.4207 & 0.2230 \\
Lending Club & TabNet & 0.6698 (0.0034) & 0.3344 & 0.4205 & 0.1612 \\
Lending Club & XGBoost & 0.6749 (0.0006) & 0.3439 & 0.4199 & 0.1600 \\
\hline
\end{tabular}
\end{table}

\subsection{Threshold selection versus resampling}
The controlled ablation shows that threshold choice materially changed default detection under no resampling. Moving from a fixed 0.50 threshold to a validation-F1 threshold increased mean F1 from 0.4758 to 0.5491 on UCI, from 0.4362 to 0.5786 on South German, and from 0.0890 to 0.4215 on Lending Club. Mean recall rose from 0.3649 to 0.5637, from 0.3750 to 0.7150, and from 0.0507 to 0.7723, respectively. Since both thresholds were applied to the same fitted model's scores, AUROC did not change; the gains reflect a different operating point, not better probability ranking. Higher recall also changes false-positive burden, so the tuned threshold should be interpreted against the stated error costs.

Resampling did not give a consistent additional gain. Under the validation-F1 threshold, no resampling had the highest mean F1 among UCI strategies (0.5491); the strongest sampler, SVM-SMOTE, averaged 0.5429. On Lending Club, no resampling averaged 0.4215 versus 0.4115 for the strongest sampler, SVM-SMOTE. On South German, ADASYN reached mean F1 0.5923 versus 0.5786 without resampling, a modest sampler advantage; validation-selected strategies also varied across seeds (SVM-SMOTE in five, SMOTE in three, and no resampling in two). By contrast, validation-AUROC selected no resampling in all ten UCI and Lending Club seeds. In the controlled ablation, no-resampling AUROC was 0.7878 on UCI and 0.6753 on Lending Club, compared with 0.7806 for the strongest UCI sampler and 0.6659 for the strongest Lending Club sampler; South German no-resampling and SMOTE both averaged 0.7843. These results support a conditional finding: threshold tuning was consistently useful, while sampler value depended on dataset and metric.

\begin{table}[ht]
\centering\small
\caption{No-resampling controlled ablation: fixed versus validation-F1 threshold. Values are held-out means across ten seeds.}
\label{tab:threshold-ablation}
\begin{tabular}{lrrrr}
\hline
Dataset & F1 (0.50) & F1 (tuned) & Recall (0.50) & Recall (tuned) \\
\hline
UCI & 0.4758 & 0.5491 & 0.3649 & 0.5637 \\
South German & 0.4362 & 0.5786 & 0.3750 & 0.7150 \\
Lending Club & 0.0890 & 0.4215 & 0.0507 & 0.7723 \\
\hline
\end{tabular}
\end{table}

\subsection{Calibration and probability quality}
The main model summaries show that probability quality did not simply follow AUROC. LightGBM had the lowest mean UCI Brier score (0.1325); Random Forest had the lowest South German Brier score (0.1684); and LightGBM (0.1599) and CatBoost (0.1603) were similar on Lending Club, while Logistic Regression was substantially worse there (0.2353). In the controlled no-resampling ablation, fitting isotonic calibration on validation predictions did not lower held-out mean Brier score: it changed 0.1325 to 0.1331 on UCI, 0.1710 to 0.1744 on South German, and 0.1602 to 0.1633 on Lending Club. For oversampled models, isotonic calibration often reduced the Brier penalty associated with resampling; for example, UCI SMOTE changed from 0.1629 uncalibrated to 0.1356 after validation-fitted isotonic calibration, and Lending Club SMOTE from 0.2178 to 0.1667. Thus calibration helped repair sampler-induced probability distortion in these examples but did not improve every no-resampling model. Calibration estimates remain split- and prevalence-dependent and should be read with the reliability bins, intercept, and slope rather than Brier score alone.

\subsection{Robustness under perturbation}
Robustness summaries average the three perturbation levels over six models and ten seeds (180 model--seed--level observations per family and dataset). Mean AUROC retention stayed above 0.95 for every family in each dataset, but threshold decisions and probability quality were more sensitive in some settings. On UCI, correlated missingness produced the lowest mean AUROC retention (0.9576); its mean Brier retention was 0.9720 and mean decision-flip rate 0.0320. On South German, numeric noise was the clearest stressor: mean AUROC retention was 0.9800, mean Brier retention 0.8362, mean absolute probability shift 0.0991, and mean decision-flip rate 0.1646. On Lending Club, mean AUROC retention was at least 0.9801 across families; missingness produced the highest mean flip rate (0.0759), while covariate shift produced 0.0105. Stable ranking therefore did not guarantee stable probabilities or unchanged classifications, particularly under South German numeric noise.

\subsection{Explanation quality and agreement}
Explanation diagnostics cover 100 test instances per seed and the validation-selected champion, giving 1,000 instance-level records per dataset. They do not compare all six classifiers, and the champion model family can vary by seed. Mean absolute SHAP--LIME rank correlation was 0.5517 on UCI, 0.4673 on South German, and 0.6739 on Lending Club, indicating moderate and dataset-dependent agreement. Mean LIME local $R^2$ was 0.2401, 0.4116, and 0.2956, respectively. Across three repeated LIME runs for the same instance and fitted model, mean rank stability was 0.5757, 0.8562, and 0.7492; mean top-10 Jaccard overlap was 0.5342, 0.7315, and 0.7626. Mean SHAP features required to cover 80\% of attribution mass were 10.9, 6.8, and 5.3. These diagnostics show that agreement, local fit, and repeatability differ by dataset. They do not measure explanation stability across refitted models, establish attribution truth, or show that an explanation is sufficient for a regulated credit decision. Comprehensiveness and sufficiency results should be reported as perturbation-based fidelity diagnostics with their training-median reference baseline, not as proof of causal validity.

\subsection{Fairness--utility trade-offs}
Fairness--utility curves and validation-fitted equal-opportunity post-processing were available for UCI and South German, using the recorded group attributes; no corresponding group analysis was possible for Lending Club. Across the evaluated UCI model--seed--attribute policies (180 comparisons), post-processing changed mean test utility by +0.0778 and accuracy by $-0.1909$. Mean TPR gap fell by 0.0335, but mean FPR gap rose by 0.1045 and mean maximum gap rose by 0.0207. For South German, 126 of 180 model--seed--attribute combinations were evaluable; mean test utility changed by +0.0125 and accuracy by $-0.0837$, while mean TPR gap rose by 0.0334, FPR gap by 0.1313, and maximum gap by 0.0923. These are unweighted descriptive averages across different attributes, models, and seeds, not population-level estimates. Validation targeting of TPR parity did not reliably reduce all held-out gaps, and utility gains came with accuracy loss. Report attribute-specific paired bootstrap intervals and per-group counts alongside curves; do not summarize this post-processor as an across-the-board fairness improvement.

\subsection{Reliability-score diagnostics}
The validation-time reliability screen combined AUROC, Brier skill, robustness retention, and inverse group gap, with minimum robustness 0.90 and maximum group gap 0.10. No candidate met both constraints in any of the ten runs for any dataset; the reported selected entry was therefore an infeasible fallback, not a feasible recommendation. A separate test-set composite score was available for all ten UCI champions (mean 0.5677), for five South German champions (the other five lacked a usable group-fairness component), and for none of the Lending Club champions. The score was calculated for the selected champion only, uses exploratory equal-weight geometric aggregation, and is not a common ranking across models or datasets. We therefore retain component-level results as findings and treat the score as a diagnostic, not as evidence of a reliable deployment choice.

The saved figure files are organized by seed rather than aggregated across all ten seeds. Any figure shown beside Table~\ref{tab:all-seed-models} must therefore be labelled as seed-specific or regenerated from all seeds. In particular, do not use the current Figure 9 test-set optimum marker as evidence for a validation-selected threshold; the plot code must be corrected or that marker omitted before publication.

\section{Discussion}
\textit{Target length: about 1,500 words.}

\subsection{Answer to the primary question}
Interpret threshold tuning as a decision-policy intervention distinct from data-level rebalancing. The main practical inference is conditional: in UCI and the packaged temporal Lending Club evaluation, no-resampling with validation-tuned thresholds performed competitively or better on the reported measures; South German indicates that resampling can still help in some settings. Do not generalize to all lenders, populations, or deployment periods.

\subsection{Relation to prior credit-scoring studies}
Compare model rankings, imbalance strategies, evaluation splits, and metrics with verified published studies. Explain which differences are explained by data, target, split, tuning budget, or threshold policy. Avoid claims of statistical superiority unless the tested comparison and multiplicity correction support them.

\subsection{Implications for model evaluation}
Recommend that credit-default benchmarks report (i) ranking and probability metrics, (ii) validation-selected operating thresholds and associated costs, (iii) paired sampler ablations, and (iv) split/seed variability. Frame this as a transparent reporting recommendation derived from this experiment, not as a newly established standard.

\subsection{Interpretability, robustness, and fairness together}
Discuss what the diagnostics add beyond predictive performance, including explanation instability, perturbation sensitivity, and group-level error differences. Explain why no single composite score substitutes for component-level inspection, stakeholder judgment, or legal review. Describe what the present score can and cannot support.

\subsection{Limitations}
Address: three datasets with different sizes and provenance; one packaged Lending Club subsample; row rather than month-grouped temporal boundaries; random splits for two datasets; seed variation changing both split and training; limited tuning budget; incomplete explanation stability analysis across refits; group fairness unavailable for Lending Club and limited to recorded attributes elsewhere; robustness perturbations not covering adversarial poisoning or broad real-world shift; exploratory composite-score design; and any data/version/hash gaps. Separate observed limitations from future work.

\subsection{Reproducibility and reporting corrections}
Ensure all published aggregate tables/figures correspond to the same set of seeds and dataset. Label seed-specific visualizations. Fix the Figure 9 marker to use the validation-selected threshold; the current plotting code recomputes the F1 optimum on the test labels. Replace absolute machine-specific manifest paths and record the checksum for the bundled Lending Club sample.

\section{Conclusion}
\textit{Target length: about 650 words.}

Answer the primary question in a few sentences. State the evidence-bounded conclusion: the operating threshold had a larger practical effect on minority-class F1 than resampling in UCI and Lending Club in these runs, whereas South German showed a more mixed sampler result. Reiterate that model rankings and resampling benefits vary across datasets. Close with the recommendation to evaluate thresholds, resampling, calibration, robustness, and fairness separately before considering any composite audit score.

% Required declarations: complete with verified author/project information before
% submission. Do not leave these as generic claims in the final manuscript.
\section*{Declarations to complete}
\begin{itemize}
  \item Data availability: identify public sources, packaged sample provenance, and any access restrictions.
  \item Ethics approval and consent: state applicability for secondary analysis of the listed datasets.
  \item Author contributions (CRediT): assign only contributions confirmed by the authors.
  \item Funding, conflicts of interest, and acknowledgements: provide verified statements.
\end{itemize}

% Bibliography is intentionally omitted until verified references and a .bib file
% are prepared. The class option selects the Springer MathPhys numeric style.
% Add \bibliography{references} when paper/references.bib exists.

\end{document}
